\chapter{Background}

\section{The memory hierarchy and load use latency}

A modern core hides memory latency with out of order execution: while one load
is outstanding, independent work continues. This makes ordinary array traversal
a poor latency probe, because the hardware overlaps many accesses. The honest
probe is a dependent load chain, where the address of the next load is the value
returned by the previous one. The processor cannot start load $n+1$ until load
$n$ retires, so the measured time per load is the true load use latency of
whichever cache level currently holds the data. Drepper~\cite{drepper2007} and
Hennessy and Patterson~\cite{hennessy2019} describe this technique and the
plateau structure it reveals.

Two things must be defeated for the probe to be honest. The stride prefetcher
will predict a regular access pattern and hide the latency, so the chase visits
a random permutation at a 64 byte (one cache line) stride, which the prefetcher
cannot anticipate. The permutation must also form a single cycle over all nodes;
a permutation that splits into shorter cycles would keep a small footprint
resident in L1 regardless of the allocation size, silently reporting L1 latency
everywhere.

\section{The STREAM model and machine balance}

McCalpin's STREAM~\cite{mccalpin1995} measures sustainable bandwidth with four
kernels, Copy, Scale, Add, and Triad, on arrays sized several times larger than
the last level cache so the data cannot hide in cache. Bandwidth is bytes moved
divided by time, where bytes moved uses a fixed convention per kernel. Non
temporal stores matter here: an ordinary store first reads the target line for
ownership, doubling write traffic, whereas a streaming store writes straight to
memory, so the counted bytes match the real DRAM traffic.

\section{The roofline model}

The roofline~\cite{williams2009} bounds attainable performance by the minimum of
two ceilings: peak compute, and arithmetic intensity multiplied by peak
bandwidth. Arithmetic intensity is FLOPs performed per byte of memory traffic. A
kernel with low intensity is memory bound and sits on the sloped bandwidth line;
a kernel with high intensity is compute bound and sits under the flat compute
ceiling. The two lines meet at the ridge point.

\section{BLIS blocking}

The BLIS framework~\cite{vanzee2015} structures matrix multiply as packed panels
that are blocked to reside at specific cache levels. The analytical
model~\cite{low2016} sizes the cache blocks, $M_c$, $K_c$, and $N_c$, from the
cache capacities and associativities so that a panel of $A$ stays in L2 and a
panel of $B$ stays in L3 while a micro panel of $B$ stays in L1. This report
computes such a prediction from the measured cache geometry and tests it against
an empirical tile sweep.
